% Part: normal-modal-logic
% Chapter: completeness
% Section: lindenbaums-lemma

\documentclass[../../../include/open-logic-section]{subfiles}

\begin{document}

\olfileid[tr]{nml}{com}{lin}

\olsection{Lindenbaum Lemması}

Lindenbaum Lemması, her $\Sigma$-tutarlı !!{formula}s kümesinin en az bir
\emph{tam} $\Sigma$-tutarlı kümenin içinde bulunduğunu kurar. Kanonik modeli
kurarken her tam $\Sigma$-tutarlı $\Delta$ kümesi için, tam olarak
$\Delta$ içindeki !!{formula}s{}in doğru olduğu bir dünya bulunduğunu
göstereceğiz. Böylece Lindenbaum Lemması her $\Sigma$-tutarlı kümenin
kanonik modeldeki bir dünyada doğru olmasını güvence altına alır.

\begin{thm}[Lindenbaum Lemması]\ollabel{thm:lindenbaum}
  $\Gamma$ $\Sigma$-tutarlıysa, $\Gamma$'yı genişleten tam
  $\Sigma$-tutarlı bir $\Delta$ kümesi vardır.
\end{thm}

\begin{proof}
Dilin bütün formüllerinin eksiksiz bir listesi $!A_0,!A_1,\dots$ olsun
(tekrarlara izin verilir). Örneğin listeye $\Obj p_0$ ile başlayıp her
$n\ge 1$ aşamasında yalnızca $\Obj p_0,\dots,\Obj p_n$ değişkenlerini
kullanan ve uzunluğu $n$ olan sonlu sayıdaki formülü listeleyin. !!{formula}s
kümeleri $\Delta_n$'yi $n$ üzerinde tümevarımla tanımlayıp
$\Delta=\bigcup_n\Delta_n$ koyacağız. Önce $\Delta_0=\Gamma$ koyalım.
$\Delta_n$ tanımlanmışsa $\Delta_{n+1}$'i şöyle tanımlayalım:
\[
\Delta_{n+1} =
\begin{cases}
  \Delta_n \cup \{!A_n\}, & \text{$\Delta_n \cup \{ !A_n\}$
    $\Sigma$-tutarlıysa;} \\
  \Delta_n \cup \{ \lnot !A_n\}, & \text{aksi durumda.}
\end{cases}
\]
Şimdi $\Delta=\bigcup_{n=0}^\infty\Delta_n$ olsun.

Bu tanımın istenen özelliklere sahip bir $\Delta$ kümesi verdiğini, yani
$\Gamma\subseteq\Delta$ ve $\Delta$'nın tam $\Sigma$-tutarlı olduğunu
göstermeliyiz.

Yapı gereği $\Delta_0\subseteq\Delta$ ve $\Delta_0=\Gamma$ olduğundan
$\Gamma\subseteq\Delta$ olduğu açıktır. Aslında her $n$ için
$\Delta_n\subseteq\Delta$; çünkü $\Delta$ bütün $\Delta_n$'lerin
birleşimidir. (Yapının her adımında önceden kurulmuş kümeye !!a{formula}
eklediğimizden $\Delta_n\subseteq\Delta_{n+1}$; $\subseteq$ geçişli
olduğundan $n\le m$ iken $\Delta_n\subseteq\Delta_m$.) Yapının her
aşamasında ya $!A_n$'yi ya $\lnot !A_n$'yi ekleriz ve her !!{formula}, bütün
$!A_n$'lerin listesinde en az bir kez görünür. Dolayısıyla her $!A$ için ya
$!A\in\Delta$ ya da $\lnot !A\in\Delta$; böylece $\Delta$ tanım gereği
tamdır.

Son olarak $\Delta$'nın $\Sigma$-tutarlı olduğunu göstermeliyiz. Bunun için
(a) $\Delta$ $\Sigma$-tutarsız olsaydı bazı $\Delta_n$'lerin
$\Sigma$-tutarsız olacağını ve (b) bütün $\Delta_n$'lerin
$\Sigma$-tutarlı olduğunu göstereceğiz.

$\Delta$'nın $\Sigma$-tutarsız olduğunu varsayalım. O zaman
$\Delta\Proves[\Sigma]\lfalse$; yani $!A_1,\dots,!A_k\in\Delta$ bulunur ve
$\Sigma\Proves !A_1\lif(!A_2\lif\cdots(!A_k\lif\lfalse)\dots)$ olur.
$\Delta=\bigcup_{n=0}^\infty\Delta_n$ olduğundan her
$!A_i\in\Delta_{n_i}$ olacak bir $n_i$ vardır. Bunların en büyüğü $n$ olsun.
$n_i\le n$ olduğundan $\Delta_{n_i}\subseteq\Delta_n$. Dolayısıyla bütün
$!A_i$'ler bir $\Delta_n$ içindedir. Bu, $\Delta_n\Proves[\Sigma]\lfalse$,
yani $\Delta_n$'nin $\Sigma$-tutarsız olduğu anlamına gelirdi.

Her $\Delta_n$'nin $\Sigma$-tutarlı olduğunu $n$ üzerinde basit
tümevarımla gösterelim. $\Delta_0=\Gamma$ ve $\Gamma$'nın
$\Sigma$-tutarlı olduğunu varsaydık; dolayısıyla iddia $n=0$ için geçerlidir.
$\Delta_n$'nin $\Sigma$-tutarlı olduğunu varsayalım. $\Delta_{n+1}$,
$\Delta_n\cup\{!A_n\}$ $\Sigma$-tutarlıysa bu küme, aksi durumda
$\Delta_n\cup\{\lnot!A_n\}$'dir. İlk durumda $\Delta_{n+1}$ açıkça
$\Sigma$-tutarlıdır. Öte yandan
\olref[prf][con]{prop:consistencyfacts}\olref[prf][con]{prop:consistencyfacts-c}
gereğince $\Delta_n\cup\{!A_n\}$ ya da
$\Delta_n\cup\{\lnot!A_n\}$ kümelerinden biri tutarlıdır; dolayısıyla öteki
durumda da $\Delta_{n+1}$ tutarlıdır.
\end{proof}

\begin{cor}\ollabel{cor:provability-characterization}
  $\Gamma\Proves[\Sigma]!A$ olması, $\Gamma$'yı genişleten her tam
  $\Sigma$-tutarlı $\Delta$ kümesi için $!A\in\Delta$ olmasına denktir.
  (Buna $\Gamma=\emptyset$ durumu da dahildir; bu durumda modal sistem
  $\Sigma$'nın başka bir nitelemesini elde ederiz.)
\end{cor}

\begin{proof}
  $\Gamma\Proves[\Sigma]!A$ olsun ve $\Delta$, $\Gamma$'yı genişleten
  gelişigüzel bir tam $\Sigma$-tutarlı küme olsun. $!A\notin\Delta$ ise
  tamlık gereğince $\lnot!A\in\Delta$; tekdüzelik gereğince
  $\Delta\Proves[\Sigma]!A$, yansımalılık gereğince de
  $\Delta\Proves[\Sigma]\lnot!A$ olur ve $\Delta$ tutarsız çıkar. Tersine
  $\Gamma\Proves/[\Sigma]!A$ ise $\Gamma\cup\{\lnot!A\}$
  $\Sigma$-tutarlıdır. Lindenbaum Lemması gereğince
  $\Gamma\cup\{\lnot!A\}$'yı genişleten tam $\Sigma$-tutarlı bir $\Delta$ kümesi
  vardır. Tutarlılık gereğince $!A\notin\Delta$.
\end{proof}

\end{document}
